\begin{table}[htbp]
\centering
\caption{Effect of Losing ZRR Status on Alternative Electoral Outcomes}
\label{tab:alt_outcomes}
\begin{tabular}{lccc}
\toprule
 & (1) & (2) & (3) \\
 & Turnout (\%) & Abstention (\%) & FN/RN Votes (count) \\
\midrule
Loser $\times$ Post & -0.041 & 0.041 & 5.483*** \\
 & (0.086) & (0.086) & (1.022) \\
 & & & \\
Commune FE & Yes & Yes & Yes \\
Year FE & Yes & Yes & Yes \\
\midrule
Observations & 72{,}387 & 72{,}387 & 72{,}387 \\
Within $R^2$ & 0.0000 & 0.0000 & 0.0012 \\
\bottomrule
\end{tabular}
\vspace{0.3em}
\parbox{\textwidth}{\footnotesize
\textit{Notes:} Difference-in-differences estimates of the effect of losing ZRR status on alternative electoral outcomes.
Column (1) uses voter turnout (\%) as the dependent variable.
Column (2) uses the abstention rate (\%), mechanically equal to $100 -$ turnout.
Column (3) uses the raw count of FN/RN votes as the dependent variable, testing whether the effect operates on the extensive margin (number of votes) rather than the intensive margin (vote share).
Treatment and control groups are identical to Table~\ref{tab:main_did}.
The observation count is slightly higher than in Table~\ref{tab:main_did} because 11 commune-elections with zero valid votes (undefined vote shares) are excluded from vote-share regressions but retained here.
All specifications include commune and year fixed effects. Standard errors clustered at the commune level in parentheses.
$^{***}p<0.01$, $^{**}p<0.05$, $^{*}p<0.1$.}
\end{table}
